% ============================================================
% OP3_1a_bridge_recon_v2.tex  (PROPOSAL-ONLY; v2 per GPT review)
% Changes vs v1: the incorrect "linear profile is no longer a
% solution" sentence REMOVED and replaced by GPT's corrected
% statement; the three notions (local EOM / stationary-density
% selection of u_*^2 / fluctuation-tensor matching) explicitly
% separated with the EOM written out; full quadratic algebra
% written out; M-(alpha,beta) kept as a matching question;
% OP3.1a phrased as completion/testing of the recorded
% OP-grad/P(X) route; E2 deferred with GPT's sentence.
% NO preprint edits.
% ============================================================
\section*{OP3.1a bridge reconnaissance, v2 (proposal-only)}

\subsection*{0. Scope and phrasing}
OP3.1a is the OP-grad completion problem: the recorded $P(X)$
motivation must be completed into a genuine bridge from bare~P3 to
the postulated~P4 ordered-gradient sector.
The recorded route is motivation, not solution; completion remains
open, and nothing here derives the mechanism.

\subsection*{1. What the article records (unchanged from v1)}
eq:PX\_action ($O(4)$-invariant $P(X)$ extension);
eq:PX\_minimum ($P_X(u_*^2)=0$, $P_{XX}(u_*^2)>0$);
eq:gradient\_branch, eq:gradient\_vac (linear ordered background,
$S^3$ degeneracy, spontaneous $O(4)\to O(3)$); status: motivation
for P4, not derived; open: form of $P$, value of $u_*$,
$u_*\leftrightarrow(\mu,\lambda)$.
Adjacent: broken-phase tensor $K^{AB}=\beta\delta^{AB}-\alpha
n^An^B$, $\alpha>\beta>0$; saddle-driven-restructuring language.

\subsection*{2. Three distinct notions (the v1 conflation,
corrected)}
For $S=\int d^4X\,P(X)$, $X=\partial_A\Phi\,\partial_A\Phi$, the
local field equation is
\[
\partial_A\bigl(P_X(X)\,\partial_A\Phi\bigr)=0 .
\]
\emph{(i) Local field equation.}
A constant-gradient profile $\partial_A\Phi=u_A$ has
$X=\mathrm{const}$, $P_X=\mathrm{const}$, and therefore solves the
local pure-$P$ field equation for \emph{any} constant value of
$P_X$, not only $P_X=0$.
\emph{(ii) Stationary-density selection of $u_*^2$.}
Among homogeneous constant-gradient configurations, the Euclidean
action density is $P(u^2)$; the recorded condition
$P_X(u_*^2)=0$, $P_{XX}(u_*^2)>0$ selects $u_*^2$ as a stationary
(minimal) action density --- a statement about magnitude
selection, distinct from~(i).
\emph{(iii) Fluctuation-tensor matching.}
A separate question again, computed next.

\subsection*{3. Quadratic algebra (written out)}
Around $\bar\Phi$ with $\partial_A\bar\Phi=u_*n_A$
($n_An_A=1$), write $\Phi=\bar\Phi+\chi$.
Then
\[
X=X_0+\delta X,\qquad X_0=u_*^2,\qquad
\delta X=2u_*\,n\!\cdot\!\partial\chi+(\partial\chi)^2 ,
\]
and
\[
P(X)=P(X_0)+P_X(X_0)\,\delta X+\tfrac12P_{XX}(X_0)\,\delta X^2
+\dots
\]
Collecting the quadratic-in-$\chi$ pieces ($\delta X$ contributes
$(\partial\chi)^2$ at first order and $4u_*^2(n\!\cdot\!\partial
\chi)^2$ from $\delta X^2$ at second order; cubic and quartic
pieces dropped):
\[
S^{(2)}\;=\;\int d^4X\Bigl[\,P_X(X_0)\,(\partial\chi)^2
+2P_{XX}(X_0)\,u_*^2\,(n\!\cdot\!\partial\chi)^2\Bigr]
\;+\;\dots
\]
Hence the effective fluctuation tensor of the pure-$P$ sector is
\[
\mathcal K^{AB}\;\propto\;P_X(X_0)\,\delta^{AB}
+2P_{XX}(X_0)\,u_*^2\,n^An^B ,
\]
with longitudinal coefficient $P_X+2P_{XX}u_*^2$ and transverse
coefficient $P_X$.
\emph{At the recorded stationary-density point} ($P_X=0$,
$P_{XX}>0$): the transverse sector receives no kinetic term from
pure $P$ at this order (leading-order flatness of the would-be
Goldstone directions of the $S^3$ degeneracy), and the
longitudinal (amplitude) sector has positive Euclidean stiffness
$2P_{XX}u_*^2$ --- consistent with $u_*$ being a density minimum.

\subsection*{4. Matching question M-($\alpha,\beta$) (corrected
statement)}
Formally matching $\mathcal K^{AB}$ to
$\beta\delta^{AB}-\alpha n^An^B$ gives
$\beta\leftrightarrow P_X(X_0)$ and
$\alpha\leftrightarrow-2P_{XX}(X_0)u_*^2$.
Matching the form $\beta\delta^{AB}-\alpha n^An^B$ with
$\alpha>\beta>0$ would require $P_X(X_0)>0$ and
$P_{XX}(X_0)<0$ at the background point.
This is not the same as the recorded pure-$P$ stationary-density
condition $P_X(u_*^2)=0$, $P_{XX}(u_*^2)>0$.
A constant-gradient profile may still solve the local pure-$P$
field equation, but the mechanism selecting its magnitude and
matching its fluctuation tensor to the broken-phase EFT tensor is
then not supplied by the recorded $P_X=0$ minimum alone.
\emph{Matching question M-($\alpha,\beta$).}
Exhibit the bridge between the recorded $P(X)$-stationary
motivation and the recorded broken-phase tensor with
$\alpha>\beta>0$: candidates include amplitude--potential coupling
along the profile (the $V(\Phi)$ sector), mixed $P(X,\Phi)$
dependence, the orientation-sector stiffness
$\mathcal C_nu^2(\partial n)^2$ entering at EFT level (natural
supplier of the transverse kinetic term left flat by pure $P$), or
a level distinction between the bare-fluctuation tensor and the
EFT tensor of the emergent-metric sector.
No inconsistency is asserted: the two tensors may legitimately
live at different levels; the task is to exhibit the bridge
explicitly.
Any resolution must remain consistent with the recorded saddle
character of the ordering configuration.

\subsection*{5. Admissible extension class (draft definition;
proposal-only)}
Single real $\Phi$; local; $O(4)$ and $\mathbb Z_2$; first pass
$P(X,\Phi)$ with $X\ge0$; Euclidean boundedness below; reduction
to eq:ECT\_action in the small-$X$/weak-coupling regime; the class
must reproduce the recorded ordered-branch EFT interface
($Z_u$, $\mathcal C_n$, $W(u)$) upon expansion, and any candidate
background must pass (i)--(iii) of \S2--\S4 separately.
Optional second pass: $(\partial^2\Phi)^2$-type terms.

\subsection*{6. Deferred routes}
E2 remains deferred; it may be useful only if the
$P(X)$-completion route fails or requires higher-gradient support.
E3 (boundary data) and E4 (measure selection) unchanged from v1;
E4 belongs with OP3.1c.

\subsection*{7. What is NOT proposed / risks}
No insertion; no status changes; no claim that the recorded
motivation is wrong; no $u_*(\mu,\lambda)$ numbers; no literature
before page-level verification.
Forbidden: ``ECT derives gradient condensation''; presenting
M-($\alpha,\beta$) as a contradiction; minimisation language for
the ordering configuration.
Safe: ``completion/testing of the recorded OP-grad/$P(X)$ route'';
``matching question recorded''.

\subsection*{8. Questions for GPT}
(Q1)~Algebra of \S3 and the corrected \S2/\S4 statements: approved
as the verified basis?
(Q2)~Eventual insertion target (after your sign-off): a short
block recording (a)~the three-notion distinction of \S2, (b)~the
fluctuation tensor of \S3 as a small computed lemma, and (c)~the
matching question M-($\alpha,\beta$) --- placed next to the
recorded $P(X)$ motivation, status ``OP-grad/OP3.1a sharpened;
completion open''?
(Q3)~Include the Goldstone-flatness reading (transverse sector
left flat by pure $P$; $\mathcal C_n$ as natural supplier) in that
block, or proposal-only?
(Q4)~After this: D1-style exact blocks for review, no edits before
that review --- agreed?
